\documentclass[11pt,a4paper]{article}

% Layout
\usepackage[margin=1.25in]{geometry}
\usepackage[utf8]{inputenc}
\usepackage[T1]{fontenc}
% Math
\usepackage{amsfonts}
\usepackage{amsmath}
\usepackage{amssymb}
\usepackage{nicefrac}
% Typography
% References and links
\usepackage[numbers,sort&compress]{natbib}
\usepackage[colorlinks=true,linkcolor=blue!60!black,citecolor=blue!60!black,urlcolor=blue!60!black]{hyperref}
\usepackage{url}
% Misc
\usepackage{booktabs}
\usepackage{xcolor}
\usepackage{graphicx}
\usepackage{parskip}
\setlength{\parskip}{6pt}
\setlength{\parindent}{0pt}

\title{\textbf{The Concept Assembly Zone}\\
\large A Dynamical Systems Framework for Cross-Layer\\
Semantic Manifold Tracking in Transformers}

\author{
  James Henry \\
  \textit{Independent Researcher} \\
  \texttt{jamesrahenry@henrynet.ca}
}

\date{March 2026}

\begin{document}

\maketitle
\thispagestyle{empty}

\begin{abstract}
	Mechanistic interpretability methods commonly extract concept representations by
	identifying the single optimal layer of a Transformer's residual stream where class
	separation peaks. This ``best layer'' heuristic is computationally efficient and
	empirically grounded, but it captures a snapshot of a process rather than the process
	itself. We introduce the \textbf{Concept Assembly Zone} (CAZ): the contiguous sequence
	of middle-to-late layers across which a semantic concept transitions from vague syntactic
	probability to a rigid, geometrically extractable direction. We formalize the CAZ
	through three layer-wise metrics---Separation, Concept Coherence, and Concept
	Velocity---and derive a principled method for identifying CAZ boundaries without
	manual layer sweeps. The framework generates four testable predictions, including
	the Mid-Stream Ablation Hypothesis: that orthogonal projection ablation applied
	\emph{within} the CAZ produces superior behavioral suppression with lower collateral
	capability damage than post-CAZ ablation, offering a candidate explanation for the
	elevated KL divergences observed in abliteration work. We situate the CAZ within the
	emerging geometric interpretability program as a complement to existing feature-extraction
	and manifold-level methods.
\end{abstract}

%-----------------------------------------------------------------------
\section{Introduction}
\label{sec:intro}

The dominant paradigm in mechanistic interpretability extracts concept representations
by identifying the single ``best layer''---the residual stream depth at which a
linear probe or difference-of-means (DoM) vector achieves maximum class
separation~\citep{zou2023representation, arditi2024refusal}. This heuristic is
computationally convenient and empirically grounded. It is also, by design, a
snapshot: it identifies the peak of a process rather than characterizing the process
itself.

Transformers are iterative dynamical systems. Each layer applies a sequence of
attention and MLP operations that \emph{write} new information into the residual
stream, modifying and extending what prior layers contributed~\citep{elhage2021mathematical}.
A concept observed at Layer~15 was shaped by Layers~10 through~14 before it;
the best-layer heuristic tells us where the concept is most legible, not how it
arrived there.

This paper introduces the \textbf{Concept Assembly Zone} (CAZ) framework, which
extends the interpretability toolkit from anatomy---\emph{where is the concept most
	visible?}---to dynamical flow---\emph{how does the concept form?} The CAZ is defined
as the specific sequence of layers where a semantic concept transitions from vague
syntactic entanglement to a rigid, mathematically extractable geometric direction.

The framework has immediate practical implications. If concepts are assembled across
a window of layers rather than peaking at a point, then:
\begin{enumerate}
	\item CAZ-windowed extraction methods may capture information present in the assembly
	      dynamics that single-layer methods do not;
	\item Ablation interventions applied during concept assembly may produce qualitatively
	      different effects than interventions applied after assembly is complete;
	\item The ``dark matter'' of unexplained model behavior~\citep{engels2024dark} may
	      partially correspond to in-progress concept construction that hasn't yet resolved
	      into the linear directions that sparse autoencoders (SAEs) are designed to capture.
\end{enumerate}

The CAZ is a theoretical framework generating specific, falsifiable predictions.
Section~\ref{sec:empirical} reports preliminary validation on the GPT-2 family
across 8 concepts and 8 architectures; frontier-scale validation is pending compute
access. We are explicit about the assumptions the framework inherits from the
broader interpretability literature and about what the current results do and do
not establish.

%-----------------------------------------------------------------------
\section{Background}
\label{sec:background}

\subsection{The Residual Stream and Concept Representation}

The residual stream formulation~\citep{elhage2021mathematical} treats each layer's
output as an additive contribution to a shared communication channel. Attention heads
and MLPs read from and write to this stream; the final residual vector is projected
onto the unembedding matrix to produce logits. This architecture makes layer-by-layer
tracking of concept geometry natural: we can ask, at each layer $l$, how well the
current residual stream separates two contrastive classes.

\subsection{Difference-of-Means and Linear Artificial Tomography}

DoM extracts a concept direction $V_\text{concept} \in \mathbb{R}^d$ as the
normalized difference between class-conditional mean activations at the chosen
layer~\citep{zou2023representation}. Linear Artificial Tomography (LAT) uses a similar
contrastive approach. Both methods produce a single vector at a single depth --- a
precise and useful representation of where the concept is most geometrically legible.
The CAZ framework asks what additional information about concept formation might be
recoverable from the layers surrounding that peak.

\subsection{Abliteration and Intervention Depth}

Arditi et al.~\citep{arditi2024refusal} demonstrated that refusal behavior across
13 open-source models is mediated by a single direction removable via weight
orthogonalization (``abliteration''). Independent replications have observed KL
divergences between abliterated and unmodified models ranging from 3.16 to
5.71---suggesting that while behavioral suppression is effective, the intervention
also affects general model capabilities. The CAZ framework offers a candidate
explanation: if abliteration applies orthogonal projection at the best layer, which
typically lies at or after the CAZ peak, the intervention may be acting on a
representation that downstream computation has already incorporated. Intervening
earlier in the assembly process may allow later layers to adapt, potentially
preserving general capabilities while maintaining behavioral suppression.

\subsection{The Emerging Geometric Program}

Gurnee et al.~\citep{gurnee2025manifolds} demonstrated that character counts are
represented on low-dimensional curved helical manifolds in the residual stream,
with attention heads performing geometric transformations on these structures.
Engels et al.~\citep{engels2025linear} found circular multi-dimensional representations
for temporal concepts (days, months, years) that are not decomposable into independent
one-dimensional SAE features. Wollschläger et al.~\citep{wollschlager2025cones}
showed that refusal occupies multi-dimensional polyhedral concept cones with multiple
independent directions. These findings establish a growing body of evidence for rich
geometric structure in activation space, and have explicitly called for unsupervised
methods to detect and characterize it. The CAZ framework is designed to complement
this geometric program by providing a layer-indexed account of when such structures
crystallize.

%-----------------------------------------------------------------------
\section{The Concept Assembly Zone}
\label{sec:caz}

\subsection{Concept Lifecycle}

By tracking the residual stream across model depth, three phases of concept formation
are empirically distinguishable:

\paragraph{Phase 1: Pre-CAZ (Context and Syntax).}
In early layers, the residual stream primarily resolves local context, grammar, and
surface token relationships. Projecting contrastive datasets into this space produces
heavily entangled activations; the separation metric is near zero. The model has not
yet committed to a semantic trajectory.

\paragraph{Phase 2: The Assembly Zone.}
In middle layers, attention heads and MLPs iteratively write semantic features into
the residual stream. Contrastive activations diverge into distinct geometrically
coherent manifolds. A stable concept direction crystallizes, and the variance
explained by the primary component spikes. This is the CAZ.

\paragraph{Phase 3: Post-CAZ (Logit Projection).}
In late layers, the model transitions from abstract representation to concrete
next-token prediction. The residual stream is projected toward the unembedding matrix.
The clean semantic manifold often degrades or re-entangles as abstract concept
geometry gives way to vocabulary-specific structure.

\subsection{Layer-Wise Metrics}

Let $h_l^{(i)} \in \mathbb{R}^d$ be the residual stream activation at layer $l$ for
sample $i$, and let $\mathcal{A}$, $\mathcal{B}$ be contrastive classes with
conditional means $\bar{h}_\mathcal{A}^{(l)}$, $\bar{h}_\mathcal{B}^{(l)}$ and
within-class covariance matrices $\Sigma_\mathcal{A}^{(l)}$, $\Sigma_\mathcal{B}^{(l)}$.

\paragraph{Separation Metric.}
We define the separation at layer $l$ using a Fisher-normalized criterion:
\begin{equation}
	S(l) = \frac{\left\|\bar{h}_\mathcal{A}^{(l)} - \bar{h}_\mathcal{B}^{(l)}\right\|_2}
	{\sqrt{\frac{1}{2}\left(\operatorname{tr}(\Sigma_\mathcal{A}^{(l)}) +
			\operatorname{tr}(\Sigma_\mathcal{B}^{(l)})\right)}}
	\label{eq:separation}
\end{equation}

Raw centroid distance is misleading when cluster dispersion varies across layers.
Early layers tend toward diffuse, high-variance representations; normalization by
within-class spread corrects for this. Mahalanobis distance would account for full
covariance structure but is numerically unstable without regularization in
high-dimensional activation spaces. Fisher normalization provides the appropriate
tradeoff between geometric fidelity and computational feasibility for initial
experiments.

\paragraph{Concept Coherence.}
Separation alone is insufficient: two classes could exhibit identical centroid
separation while one forms a tight cluster and the other a diffuse cloud. We track
Concept Coherence as the explained variance ratio of the first principal component
of the between-class direction at each layer:
\begin{equation}
	C(l) = \frac{\lambda_1^{(l)}}{\sum_i \lambda_i^{(l)}}
	\label{eq:coherence}
\end{equation}
where $\lambda_i^{(l)}$ are the eigenvalues of the pooled activation covariance at
layer $l$, projected onto the contrastive subspace. A concept is \emph{well-formed}
when both $S(l)$ and $C(l)$ are high: the classes are far apart \emph{and} the
separating direction is geometrically clean.

\paragraph{Concept Velocity.}
To identify CAZ boundaries, we compute the rate of geometric divergence between layers.
Because raw layer-to-layer differences are noisy, we apply a smoothed estimate:
\begin{equation}
	v_\text{concept}(l) = \frac{1}{2k+1} \sum_{j=l-k}^{l+k} \left[S(j) - S(j-1)\right]
	\label{eq:velocity}
\end{equation}
where $k$ is the smoothing half-window. A practical heuristic is
$k = \lfloor L/24 \rfloor$, where $L$ is total model depth, yielding $k=1$ for
12--24 layer models, $k=2$ for 48-layer models, and $k=3$ for 72-layer models.
This scales the smoothing window proportionally to model depth and prevents false
CAZ boundary detection from single anomalous layers. The appropriate value of $k$
should ultimately be determined empirically: for models where ground-truth concept
boundaries can be established via ablation, the $k$ value that maximizes boundary
prediction accuracy is preferred.

\subsection{CAZ Boundary Detection}

The CAZ boundaries are derived from the velocity profile:
\begin{itemize}
	\item \textbf{CAZ Entry} ($l_\text{start}$): The first layer where
	      $v_\text{concept}(l)$ exceeds a sustained positive threshold $\theta_+$.
	\item \textbf{CAZ Peak} ($l_\text{max}$): The layer where $S(l)$ reaches its
	      absolute maximum. This corresponds to the ``best layer'' of conventional
	      interpretability.
	\item \textbf{CAZ Exit} ($l_\text{end}$): The layer where $v_\text{concept}(l)$
	      becomes consistently negative, marking the onset of post-CAZ degradation.
\end{itemize}

The conventional best-layer heuristic extracts $V_\text{concept}$ at $l_\text{max}$.
CAZ-aware extraction uses the full interval $[l_\text{start}, l_\text{end}]$.

\subsection{Multi-Layer Concept Extraction}

Three extraction methods should be compared empirically:

\begin{enumerate}
	\item \textbf{Delta PCA}: PCA on layer-to-layer residual deltas
	      $\Delta h_l = h_l - h_{l-1}$ within $[l_\text{start}, l_\text{end}]$.
	      This captures what each layer \emph{adds}---the construction process itself.
	\item \textbf{Windowed PCA}: PCA on raw activations $h_l$ across
	      $[l_\text{start}, l_\text{end}]$. This captures the cumulative concept
	      direction as it evolves through the assembly zone.
	\item \textbf{Single-layer (baseline)}: The standard DoM vector at $l_\text{max}$.
\end{enumerate}

If methods (1) or (2) consistently outperform (3) on downstream steering and
classification tasks, this validates the dynamical framing as more than descriptive
vocabulary. If they do not, the framework may still provide useful theoretical
structure while failing to improve practical extraction---which is itself worth
establishing.

%-----------------------------------------------------------------------
\section{Testable Predictions}
\label{sec:predictions}

The CAZ framework generates four predictions that are in principle falsifiable with
existing open-weight models and standard interpretability tooling.

\subsection{The Mid-Stream Ablation Hypothesis}

\textbf{Prediction 1}: Post-CAZ ablation is destructive; early-CAZ ablation is
graceful.

When orthogonal projection ablation is applied after the CAZ, the intervention acts
as an amputation. The model has already organized its output strategy around the
concept; removing the concept vector causes representational shock---elevated KL
divergence, increased perplexity, and degraded coherence on unrelated tasks.
Intervening \emph{within} the early CAZ should produce qualitatively different
behavior. If the concept is projected out of the residual stream as it is being
assembled, subsequent attention heads in later layers can route around the missing
information. The model should default to a neutral or safe state while preserving
general capabilities.

\textbf{Experimental design}: For a given concept (e.g., refusal), extract the
concept direction at each layer via DoM. Apply orthogonal projection at each layer
$l$ independently. Measure: (a) behavioral suppression rate on targeted prompts,
(b) KL divergence from the unmodified model on unrelated prompts, (c) perplexity
on a held-out benchmark. Plot the ratio $\nicefrac{(a)}{(b)}$---behavioral suppression
per unit of collateral damage---as a function of layer. The CAZ framework predicts
this ratio peaks within $[l_\text{start}, l_\text{end}]$ and degrades sharply
post-$l_\text{end}$.

This prediction connects naturally to the abliteration literature~\citep{arditi2024refusal}.
The CAZ framework suggests that systematically varying intervention depth --- and
measuring the suppression-to-capability-impact ratio as a function of layer --- would
provide a principled basis for selecting intervention points. The framework predicts
that early-CAZ intervention would improve this ratio relative to post-peak intervention.

\subsection{Architecture-Stable CAZ Positioning}

\textbf{Prediction 2}: CAZ boundaries are concept-specific but architecture-stable.

Different concepts (refusal, honesty, credibility) should have different CAZ windows
within the same model. However, the \emph{relative} positioning of those windows---as
a fraction of total model depth---should be consistent across architectures of similar
parameter count and training regime. Preliminary cross-architecture results showing
convergence at relative layer 12 for credibility encoding across multiple
GPT-family architectures are suggestive of this pattern, though not yet confirmed
on production-scale models.

\subsection{CAZ Width and Concept Abstraction}

\textbf{Prediction 3}: CAZ width correlates with concept abstraction level.

More abstract concepts (e.g., ``trustworthiness,'' ``moral valence'') should have
wider CAZ windows than concrete ones (e.g., ``negation,'' ``plurality''), because
abstract concepts require more iterative construction across attention layers.
This is testable by comparing $l_\text{end} - l_\text{start}$ for concepts at
different levels of semantic abstraction as operationalized by, for example,
depth in WordNet or scores on standard concreteness rating datasets.

\subsection{Post-CAZ Degradation as Logit Interference}

\textbf{Prediction 4}: Post-CAZ re-entanglement correlates with unembedding
matrix structure.

The degradation of clean concept geometry in late layers is not noise but a
structural consequence of preparing the residual stream for logit projection.
Concepts whose associated vocabulary tokens are distributionally similar in the
unembedding space---close in embedding distance---should show more post-CAZ
degradation than concepts with distributionally distinct vocabulary. This would
explain why some concepts retain clean geometry into late layers (their vocabulary
is well-separated) while others degrade early (their vocabulary clusters).

%-----------------------------------------------------------------------
\section{Relationship to Existing Work}
\label{sec:related}

\paragraph{Manifold interpretability.}
Gurnee et al.~\citep{gurnee2025manifolds} found curved manifolds in middle layers
for character counting and explicitly called for unsupervised geometric discovery
methods. The CAZ framework provides a formalism for identifying \emph{where} in the
layer stack such manifolds crystallize---the assembly zone is precisely where
curved manifold structure should be most geometrically coherent and most amenable
to unsupervised detection.

\paragraph{SAE dark matter.}
Engels et al.~\citep{engels2024dark} investigated the structured residual
(``dark matter'') left unexplained by sparse autoencoders, finding that approximately
50\% of the SAE error vector and over 90\% of its norm can be linearly predicted
from the input activation---the unexplained residual is structured, not noise---and
that larger SAEs largely fail to reconstruct the same contexts as smaller SAEs.
Crucially, they distinguish a linearly predictable error component from a genuinely
\emph{nonlinear} residual that behaves differently from the rest. The CAZ phase model
offers a candidate mechanism for part of this nonlinear component: SAEs are trained on
activations at a fixed layer, but in-progress concept construction within the CAZ
produces transitional representations that are neither fully syntactic (pre-CAZ) nor
fully semantic (post-CAZ). These transitional states may be precisely the activations
that resist linear decomposition regardless of SAE scale, connecting the assembly-zone
framing to the structured nature of the dark matter Engels et al.\ observe.

\paragraph{Abliteration.}
Arditi et al.~\citep{arditi2024refusal} and Wollschläger et al.~\citep{wollschlager2025cones}
establish the geometry of refusal. The CAZ framework extends this by asking not
just \emph{what} the geometry is but \emph{when} it forms, and using that temporal
structure to identify optimal intervention points. Wollschläger et al.'s finding
that refusal occupies multi-dimensional concept cones raises the open question of
whether cone dimensionality varies across the CAZ---whether it narrows during
assembly (becoming more extractable) and expands during post-CAZ degradation.

\paragraph{Representation engineering.}
Zou et al.~\citep{zou2023representation} extract directions for honesty, morality,
power-seeking, and related concepts via contrastive stimuli. The CAZ framework
provides a theoretical basis for choosing intervention depth rather than treating
it as a hyperparameter, and suggests that the full CAZ window may contain richer
directional information than any single layer.

%-----------------------------------------------------------------------
\section{Preliminary Empirical Results}
\label{sec:empirical}

The framework has been partially validated on the GPT-2 family using a full empirical
pipeline across 8 concepts $\times$ 8 model architectures (GPT-2, GPT-2-Medium,
GPT-2-Large, GPT-2-XL, GPT-Neo, Pythia, OPT variants), implemented in the companion
repository Rosetta Manifold.\footnote{\url{https://github.com/jamesrahenry/Rosetta_Manifold}}

\subsection{GPT-2-XL Results}

At GPT-2-XL scale (48 layers, 100 contrastive pairs per concept):

\begin{table}[h]
	\centering
	\begin{tabular}{llccc}
		\toprule
		\textbf{Concept} & \textbf{Type} & \textbf{Peak layer} & \textbf{Rel.\ depth} & \textbf{Peak $S$} \\
		\midrule
		temporal\_order  & relational    & L36 / 48            & 75\%                 & 0.449             \\
		causation        & relational    & L37 / 48            & 77\%                 & 0.488             \\
		negation         & syntactic     & L39 / 48            & 81\%                 & 0.314             \\
		certainty        & epistemic     & L44 / 48            & 92\%                 & 0.500             \\
		moral\_valence   & affective     & L44 / 48            & 92\%                 & 0.294             \\
		sentiment        & affective     & L44 / 48            & 92\%                 & 0.396             \\
		credibility      & epistemic     & L46 / 48            & 96\%                 & 0.736             \\
		plurality        & syntactic     & L47 / 48            & 98\%                 & 0.322             \\
		\bottomrule
	\end{tabular}
	\caption{CAZ peak layers and separation scores for 8 concepts at GPT-2-XL scale
		(48 layers). Relative depth is peak layer / total depth.}
	\label{tab:results}
\end{table}

\subsection{Confirmed Findings}

\textbf{Prediction~1 (Mid-Stream Ablation Hypothesis)} is confirmed at GPT-2 scale:
single-layer orthogonal projection ablation at the CAZ peak achieves 100\% separation
reduction on concept probes. This result does not extend to GPT-2-XL, where single-layer
projection is insufficient---consistent with the expectation that 1.5B-scale models
require multi-layer or weight-space intervention.

\textbf{Broad late-assembly ordering.} A clear two-cluster structure is observed.
Relational and syntactic concepts (\textit{temporal\_order}, \textit{causation},
\textit{negation}) assemble in the 75--81\% depth range; affective and epistemic
concepts (\textit{certainty}, \textit{moral\_valence}, \textit{sentiment},
\textit{credibility}) cluster in the 92--96\% range. This separation is robust across
the 8 architectures tested.

\textbf{Credibility} exhibits substantially stronger separation signal ($S = 0.736$)
than all other concepts---more than 50\% larger than the next strongest
(\textit{certainty}, $S = 0.500$)---suggesting a particularly well-defined concept
direction.

\subsection{Anomalies and Open Questions}

\textbf{Within-type ordering is reversed.} The framework predicted syntactic concepts
would assemble earlier than relational ones; the data shows the opposite---causation
and temporal\_order (relational) assemble at 75--77\%, while negation (syntactic)
assembles later at 81\%. The between-type ordering holds; the within-type prediction
does not.

\textbf{Plurality is anomalously deep} (L47, 98\% depth)---the deepest concept
measured, deeper even than credibility. A syntactic concept assembling in the final
layer of a 48-layer model is not explained by the current framework and warrants
investigation at larger scale.

\textbf{Prediction~2 (Architecture-Stable Positioning)} is partially supported: the
broad two-cluster ordering holds across architectures, but proper validation requires
same-scale cross-architecture comparison (e.g., GPT-Neo-1.3B vs.\ OPT-1.3B vs.\
GPT-2-XL at matched parameter count).

\textbf{Predictions~3 and~4} (CAZ width and post-CAZ logit interference) have not yet
been tested and require either frontier-scale models or a wider concept vocabulary with
operationalized abstraction levels.

\subsection{Frontier-Scale Validation}

Frontier-scale validation (Llama~3 70B, Qwen~2.5 72B, Mistral Large) is pending compute
access. The proxy-scale results are sufficient to establish that the CAZ is a real and
measurable phenomenon; they are not sufficient to confirm the architecture-stability and
abstraction-width predictions at production model scale.

%-----------------------------------------------------------------------
\section{Limitations}
\label{sec:limitations}

Known limitations include:

\textbf{Computational cost.} Computing $S(l)$ and $C(l)$ at every layer for large
contrastive datasets and large models is expensive. Efficient approximations (random
projection, activation sketching) may be necessary for models beyond 7B parameters.

\textbf{Smoothing sensitivity.} CAZ boundary detection depends on the smoothing
parameter $k$ and threshold $\theta_+$. A principled method for setting these---perhaps
based on layer count or activation dimensionality---is needed before the framework
can be applied systematically.

\textbf{Linearity assumption.} The separation metric assumes the concept manifold
is approximately linearly separable. For concepts with curved or multi-dimensional
structure~\citep{gurnee2025manifolds, engels2025linear}, kernel-based or topological
metrics may be required. This limitation is shared with most of the current
interpretability literature.

\textbf{Token position dependence.} The framework does not account for which token
positions carry concept information. Zhao et al.~\citep{zhao2025harmfulness} showed
that harmfulness and refusal encode at different token positions. CAZ boundaries
may vary by token position as well as by layer.

\textbf{Causal vs.\ correlational.} High separation at a layer does not establish
that the concept is \emph{used} at that layer for downstream computation. Causal
validation via ablation and activation patching is required to distinguish
geometrically present concepts from epiphenomenal structure.

%-----------------------------------------------------------------------
\section{Conclusion}
\label{sec:conclusion}

The Concept Assembly Zone provides a framework for analyzing how Transformers
construct semantic representations across their depth. By shifting from single-layer
snapshots to cross-layer flow analysis, it generates specific, testable predictions
about concept extraction quality, optimal ablation depth, and the geometric origins
of the SAE dark matter.

The framework's value is empirical, not definitional. Preliminary results on the
GPT-2 family confirm that the CAZ is a measurable phenomenon: a robust two-cluster
ordering of concept assembly depth is observed across 8 concepts and 8 architectures,
and Prediction~1 (Mid-Stream Ablation) is confirmed at GPT-2 scale. Anomalies
remain---the within-type syntactic/relational ordering is reversed, and plurality
assembles anomalously late---and Predictions~2 through~4 await frontier-scale
validation.

The next critical experiment is Prediction~1 at production scale: for a model such
as Llama~3 70B and a well-studied concept such as refusal, compare the ratio of
behavioral suppression to collateral capability damage as a function of intervention
layer. If this ratio peaks within the CAZ window, it provides both practical guidance
for safer ablation and theoretical support for the dynamical framing. If it does not,
the framework's descriptive vocabulary may still be useful while its interventional
claims are refuted. Either outcome advances the field.

The CAZ framework is offered as a complement to existing interpretability methods:
extending the layer-sweep heuristic into a dynamical account, and connecting
single-layer feature extraction to the emerging geometric program that the field
is converging toward. Whether it survives contact with data is an open question
we invite others to test.

%-----------------------------------------------------------------------
\bibliographystyle{plainnat}
\bibliography{caz_refs}

\end{document}
